%% kapa: exam=2 round=19 year=2008 period=1 subject=appraisal_practice subject_ko=감정평가실무 seq=3 number=3 points=15
%% source: Q-Net 감정평가사 기출문제 (cst00309) artlSeq=5024177 file=extracted/1008/5024177/2008년 제19회 감정평가사 2차시험 기출문제/1교시.hwp
%% conversion: kordoc 4.24.51 -> tools/gen2.py (auto); review=proofread-sonnet; flags=none
\begin{longquestion}{【문제 3】}
A감정평가사는 ○○청으로부터 아래와 같은 내용의 입목에 대한 감정평가의뢰를 받았다. 제시자료를 검토하여 입목의 취득가격을 결정하시오. (단, 입목의 평가방법은 제시자료에 타당한 합리적이고 보편적인 방식을 선택하여 평가할 것) (15점) \par
<자료 1> 감정평가 의뢰내역 \par
1) 개요 \par
① 평가목적 : 조림대부지 내 입목의 취득(매수) \par
② 소재지 : ○○○도 ○○군 ○○면 ○○리 산21 \par
③ 지목 : 임야 \par
④ 면적 : 1,050,000 m\textsuperscript{2} \par
2) 입목현황 \par
{\small\centering
\setlength{\tabcolsep}{3pt}
\begin{tabular}{|c|c|c|c|c|c|c|c|c|}
\hline
\makecell{임종} & \makecell{임상} & \makecell{수종} & \makecell{혼효율\newline{}(\%)} & \makecell{임령} & \makecell{령급} & \makecell{경급\newline{}(cm)} & \makecell{수고\newline{}(m)} & \makecell{ha당 재적(m\textsuperscript{3})} \\
\hline
\multirow{2}{*}{\makecell{천연림\newline{}(자연림)}} & \makecell{활엽수} & \makecell{참나무\newline{}기타활엽수} & \multirow{2}{*}{\makecell{70}} & \multirow{2}{*}{\makecell{$\frac{ 29 } { 15-45 }$}} & \multirow{2}{*}{\makecell{Ⅱ-Ⅴ}} & \multirow{2}{*}{\makecell{$\frac{ 18 } { 8-35 }$}} & \multirow{2}{*}{\makecell{$\frac{ 10 } { 8-18 }$}} & \multirow{2}{*}{\makecell{75}} \\
\cline{2-3}
 & \makecell{침엽수} & \makecell{소나무} &  &  &  &  &  &  \\
\hline
\makecell{인공림\newline{}(조림)} & \makecell{침엽수} & \makecell{잣나무\newline{}낙엽송\newline{}리기다소나무} & \makecell{30} & \makecell{$\frac{ 35 } { 25-45 }$} & \makecell{Ⅲ-Ⅳ} & \makecell{$\frac{ 20 } { 10-36 }$} & \makecell{$\frac{ 11 } { 8-19 }$} & \makecell{95} \\
\hline
\end{tabular}\par}

※ 참고사항 : 1. 조림대부지로서 입목의 관리상태는 양호함. \par
2. 경급(cm) : $\dfrac{\text{평균경급}}{\text{최저경급-최고경급}}$ \par
3) 수종별 재적 \par
{\small\centering
\begin{tabular}{|c|c|c|c|c|}
\hline
\makecell{임종} & \makecell{임상} & \makecell{수종} & \makecell{재적(m\textsuperscript{3})} & \makecell{비고} \\
\hline
\multirow{4}{*}{\makecell{천연림\newline{}(자연림)}} & \multirow{2}{*}{\makecell{활엽수}} & \makecell{참나무} & \makecell{1,653.80} & \multirow{4}{*}{} \\
\cline{3-4}
 &  & \makecell{기타활엽수} & \makecell{3,307.50} &  \\
\cline{2-4}
 & \makecell{침엽수} & \makecell{소나무} & \makecell{551.30} &  \\
\cline{2-4}
 & \multicolumn{2}{|c|}{\makecell{소계}} & \makecell{5,512.60} &  \\
\hline
\multirow{4}{*}{\makecell{인공림\newline{}(조림)}} & \multirow{3}{*}{\makecell{침엽수}} & \makecell{잣나무} & \makecell{1,047.40} & \multirow{4}{*}{} \\
\cline{3-4}
 &  & \makecell{낙엽송} & \makecell{748.10} &  \\
\cline{3-4}
 &  & \makecell{리기다소나무} & \makecell{1,197.00} &  \\
\cline{2-4}
 & \multicolumn{2}{|c|}{\makecell{소계}} & \makecell{2,992.50} &  \\
\hline
\multicolumn{3}{|c|}{\makecell{합계}} & \makecell{8,505.10} &  \\
\hline
\end{tabular}\par}

<자료 2> 입목 평가자료 \par
1) 원목 시장가격(가격시점현재) \par
{\small\centering
\begin{tabular}{|c|c|c|c|c|c|c|c|}
\hline
\multirow{2}{*}{\makecell{등급\newline{}기준}} & \multirow{2}{*}{\makecell{흉고직경\newline{}(경급)}} & \multicolumn{6}{|c|}{\makecell{원목가격(원/m\textsuperscript{3})}} \\
\cline{3-8}
 &  & \makecell{참나무} & \makecell{기타\newline{}활엽수} & \makecell{소나무} & \makecell{잣나무} & \makecell{낙엽송} & \makecell{리기다\newline{}소나무} \\
\hline
\makecell{상} & \makecell{30cm 이상} & \makecell{105,000} & \makecell{100,000} & \makecell{110,000} & \makecell{100,000} & \makecell{105,000} & \makecell{100,000} \\
\hline
\makecell{중} & \makecell{16cm 이상} & \makecell{90,000} & \makecell{85,000} & \makecell{95,000} & \makecell{90,000} & \makecell{95,000} & \makecell{90,000} \\
\hline
\makecell{하} & \makecell{16cm 미만} & \makecell{85,000} & \makecell{78,000} & \makecell{85,000} & \makecell{80,000} & \makecell{85,000} & \makecell{80,000} \\
\hline
\end{tabular}\par}

※ 용재림 및 기타용도(펄프, 갱목, 목탄 및 목초액의 용도 등)등으로 사용할 수 있는바 일반기준 벌기령은 적용하지 아니하고, 시장가격은 천연림과 인공림(조림)의 구분없이 형성되고 있음. \par
2) 조재율 \par
(단위 : \%) \par
{\small\centering
\begin{tabular}{|c|c|c|}
\hline
\makecell{등급기준} & \makecell{활엽수} & \makecell{침엽수} \\
\hline
\makecell{상} & \makecell{90} & \makecell{90} \\
\hline
\makecell{중} & \makecell{85} & \makecell{85} \\
\hline
\makecell{하} & \makecell{80} & \makecell{80} \\
\hline
\end{tabular}\par}

3) 생산비용 \par
① 벌목조재비 \par
{\small\centering
\begin{tabular}{|c|c|c|c|}
\hline
\multicolumn{2}{|c|}{\makecell{1일 노임/인}} & \multirow{2}{*}{\makecell{기계상각비 및 연료비}} & \multirow{2}{*}{\makecell{1일 작업량/인}} \\
\cline{1-2}
\makecell{벌목비} & \makecell{조재비} &  &  \\
\hline
\makecell{80,000원} & \makecell{80,000원} & \makecell{30,000원} & \makecell{10.0m\textsuperscript{3}} \\
\hline
\end{tabular}\par}

② 산지집재비(소운반 포함) \par
1일 노임은 80,000원/인이며 1일 작업량은 10.0m\textsuperscript{3}/인임 \par
③ 운반비 \par
{\small\centering
\begin{tabular}{|c|c|c|}
\hline
\makecell{구분} & \makecell{1일 노임/인} & \makecell{1일 작업량/인} \\
\hline
\makecell{상하차비} & \makecell{80,000원} & \makecell{10.0m\textsuperscript{3}} \\
\hline
\makecell{자동차운반비} & \makecell{110,000원} & \makecell{10.0m\textsuperscript{3}} \\
\hline
\end{tabular}\par}

④ 임도 보수 및 설치비 \par
{\small\centering
\begin{tabular}{|c|c|c|}
\hline
\makecell{1일 노임/인} & \makecell{1일 작업량/인} & \makecell{소요임도} \\
\hline
\makecell{90,000원} & \makecell{0.3km} & \makecell{2.1km} \\
\hline
\end{tabular}\par}

⑤ 잡비 : 생산비용의 10\% \par
4) 이자율 및 기업자이윤 등 \par
① 자본회수기간은 6개월정도이며 이자율은 금융기관 대출금리기준 연 7.0\%를 적용함. \par
② 기업자 이윤은 10\%, 산재보험을 포함한 위험률은 5.0\%로 적용함. \par
<자료 3> 참고사항 \par
1) 일부 수종에서 참나무 시들음병이 발생되어 피해도가 “중”이상인 입목은 평가에서 제외하고 피해도가 “경”이하인 입목은 정상입목 평가액의 90\%수준으로 평가함이 적절함. \par
2) 참나무 시들음병 피해도를 조사한 바, 조사재적 중 “중”이상 입목은 약 50\%(826.90m\textsuperscript{3}), “경”이하 입목은 약 20\%(330.80m\textsuperscript{3})임. \par
3) 단가 계산은 원단위는 절사하고 십원단위까지만 표기요함. \par
\end{longquestion}
